\subsection{INTERVALS IN SETS OF INTEGERS}

Every set of integers, being a set of cardinals, is well-ordered (§ 3, no. 2, Theorem 1). Furthermore, for each integer \(a\) the relation `` \(x\) is a cardinal and \(x \leqslant a\) '' is collectivizing in \(x\) (§ 3, no. 2, Remark after Theorem 1), and the set of \(x\) which satisfy this relation is a set of integers (§ 4, no. 2, Proposition 2), which may therefore be denoted by {[}0, a{]}.

PROPOSITION 4. Let \(a\) and \(b\) be integers. Then the mapping \(x \to a + x\) is a strictly increasing isomorphism of the interval \([0, b]\) onto the interval \([a, a + b]\) , and \(y \to y - a\) is the inverse isomorphism.

Clearly the relations \(0 \leqslant x \leqslant b\) imply \(a \leqslant a + x \leqslant a + b\) . The mapping \(x \to a + x\) is strictly increasing (and therefore injective) by Proposition 3 of no. 2. Finally, the relations \(a \leqslant y \leqslant a + b\) imply \(y = a + x\) with \(x \geqslant 0\) and \(a + x \leqslant a + b\) , whence \(x \leqslant b\) (no. 2, Proposition 3). This completes the proof.

PROPOSITION 5. If \(a\) and \(b\) are integers such that \(a \leqslant b\) , then the interval \([a, b]\) is a finite set whose number of elements is \(b - a + 1\) .

By Proposition 4 we may restrict ourselves to the case where \(a=0\) . The proof is by induction on \(b\) . If \(b=0\) , the result is clear. The relation \(0 \leqslant x \leqslant b+1\) is equivalent to `` \(0 \leqslant x < b+1\) or \(x=b+1\) '', and the relation \(0 \leqslant x < b+1\) is equivalent to \(0 \leqslant x \leqslant b\) (§ 4, no. 2, Proposition 2); in other words, the interval \([0, b+1]\) is the union of \([0, b]\) and \(\{b+1\}\) , and these two sets are disjoint. By the inductive hypothesis the number of elements of \([0, b+1]\) is equal to \((b+1)+1\) , and the Proposition is proved.

PROPOSITION 6. For every finite set \(\mathbf{E}\) which is totally ordered and has \(n\) elements \((n \geqslant 1)\) , there exists a unique isomorphism of \(\mathbf{E}\) onto the interval \([1, n]\) .

Since E and \([1, n]\) are well-ordered (§4, no. 4, Proposition 3, Corollary 1) and have the same number of elements (Proposition 5), the result follows from Theorem 3 of §2, no. 5 and Corollary 2 to Proposition 2 of §4, no. 2.

